\section{Reproducibility and source assets}
\label{app:reproduction}

\paragraph{Delivered implementation.}
The delivery root contains the report source, built \texttt{main.pdf}, analysis package, configurations and presentation scripts. The commands below reproduce the frozen main pipeline; the completed raw checkpoints and numerical summaries are supplied in this delivery. The current evidence status is recorded in \texttt{results/main\_execution\_manifest.json}. That manifest binds the executable configuration and attribution sources to their SHA-256 identity. The complete inventory comprises 540 reference games and 97,200 stochastic calls. Main metadata retain raw-data hashes, source/configuration signatures, seeds, selected input/background IDs, preprocessing groups, library versions and hardware. Checkpoints are reused only when their configuration and source identity match. The independent panel audit checks saved attribution/reference vectors and complete run denominators before aggregation.

\paragraph{Environment and reproduction commands.}
The analysis uses Python 3.11, with exact top-level packages in \texttt{analysis/requirements.txt} and installed versions in \texttt{analysis/environment-versions.json}. The analysis package specifies single-threaded BLAS/OpenMP execution and no GPU. From the delivery root, create a Python environment, install the pinned packages, then run:
\begin{verbatim}
python3 -m venv .venv
. .venv/bin/activate
python3 -m pip install -r analysis/requirements.txt
cd analysis
export OPENBLAS_NUM_THREADS=1
export OMP_NUM_THREADS=1
export MKL_NUM_THREADS=1
python3 scripts/fetch_data.py
PYTHONPATH=src python3 tests/test_estimators.py
PYTHONPATH=src python3 tests/test_metrics.py
python3 scripts/run_queue.py \
  german:logistic german:forest german:histboost \
  acs:logistic acs:forest acs:histboost \
  diabetes:logistic diabetes:forest diabetes:histboost \
  --output ../results/reproduction/main
python3 scripts/audit_panel.py \
  --input ../results/reproduction/main \
  --output ../results/reproduction/panel_audit.json
python3 scripts/aggregate.py \
  --input ../results/reproduction/main \
  --output ../results/reproduction
python3 scripts/evidence_handoff.py \
  --input ../results/reproduction
python3 scripts/audit_timing.py \
  --input ../results/reproduction
\end{verbatim}
To regenerate explanations, use a fresh output root: the delivered checkpoints omit fitted joblib bundles. The acquisition script checks the recorded dataset hashes; it does not silently substitute changed data. Separate cell queues can be run in parallel when memory permits. Actual elapsed time depends on that scheduling and on shared-host contention. The accepted pre-execution ledger projected approximately 7.4 CPU hours for the main matrix, or 13.3 CPU hours including planning allowances and reserves, against the objective's approximate 24-CPU-hour planning scale. The measured main total is 7.493 CPU hours on Intel Xeon Gold 6230 CPUs at 2.10 GHz in Linux containers with two-core affinity and one-thread numerical libraries. Reproduction should allow roughly eight CPU hours for the main matrix plus acquisition, preliminary checks, statistical postprocessing and any separately executed sensitivity arm; scheduling determines elapsed duration. The recorded original wall clocks include severe interruptions and are preserved in \texttt{results/timing\_audit.json}, rather than used as a clean runtime prediction.

To audit and reaggregate the supplied evidence directly, use the same audit/aggregate commands with input \texttt{../results/main} and outputs \texttt{../results/}, without running acquisition or model fitting. The fresh-root reproduction preflight independently regenerated one Diabetes logistic main input and all 180 calls with exactly matching saved reference/attribution vectors (\texttt{results/reproduction\_preflight.json}); it is not an additional statistical replication. For the separate background arm, after fresh main reproduction run:
\begin{verbatim}
python3 scripts/run_background_queue.py \
  --main-results ../results/reproduction \
  --output ../results/reproduction/background
python3 scripts/aggregate_background.py \
  --main-results ../results/reproduction \
  --input ../results/reproduction/background \
  --output ../results/reproduction/background_sensitivity
\end{verbatim}
Audit the delivered background checkpoints instead with main-results \texttt{../results}, input \texttt{../results/background} and output \texttt{../results/background\_sensitivity}. The shipped main and background raw vectors are already present; reaggregation requires no new explanations.

\paragraph{Whole-study computation accounting.}
\texttt{results/computation\_ledger.json} and its CSV record the known and uninstrumented executed stages, source hashes and deduplication rules. Disjoint process-clock intervals sum to 28,206.959 seconds (7.835266 CPU hours), including the 7.49345-hour main jobs, four distinct preliminary timing stages, ACS preflight, three development smokes, a fresh reproduction preflight and the background pilot/full phase timers. This is an incomplete recorded subtotal. The four preliminary stages together have 708.346 recorded CPU seconds; their startup/loading/writing outside those clocks is unknown. The unsuccessful TreeExplainer strict abort/diagnostics, early precision/parser corrections, final common-game validation and repeated local tests/audits/postprocessing retain explicit untimed entries. None is assigned zero or estimated from wall time.

Managed resource observations are separate, sampled and nonadditive: the archived cumulative counter reports 7.855328 CPU hours, and managed main/background postprocessing observations are 46.904/2.965 seconds. These overlap process clocks, omit sandbox-only work and can undershoot a phase clock locally. Neither adding the counter to the subtotal nor calling either an exact all-inclusive cost is justified. The supplied approximate 24-CPU-hour scale was a planning estimate; absent instrumentation limits the achieved total we can establish. \texttt{python3 scripts/computation\_ledger.py --input ../results}, run from \texttt{analysis/}, regenerates the ledger from preserved source records without rerunning experiments.

\paragraph{Predictor details.}
The fixed model configuration is logistic regression with $C=1$ and at most 2,000 iterations; a random forest with 100 trees, minimum leaf size 5 and a feature fraction of $.8$; and histogram gradient boosting with 100 iterations, at most 15 leaves, L2 regularization 1 and early stopping disabled. Other constructor defaults are fixed by scikit-learn 1.9.1. These configurations were selected for contrasting CPU-feasible predictors rather than tuned to explanation quality. Numeric preprocessing is fitted per raw feature using training-only median imputation and standardization. Categorical preprocessing uses training-only most-frequent imputation and either one-hot encoding for logistic regression or ordinal codes for trees; unknown categories are ignored by the former and mapped to $-1$ by the latter. The original-feature group map is retained through the transformation. Train/test splitting and model training use seed 20261103, capped training subsampling uses seed 20261104, and background/panel selection uses seed 20261105. Explainer seeds and bootstrap seeds are stated in Section~\ref{sec:study} and retained in run metadata.

\paragraph{German Credit.}
We use Hofmann's Statlog German Credit data from UCI, DOI \href{https://doi.org/10.24432/C5NC77}{10.24432/C5NC77}, distributed under \href{https://creativecommons.org/licenses/by/4.0/}{CC BY 4.0} \citep{hofmann1994}. The original symbolic \texttt{german.data} supplies the 20 feature players, rather than the separately recoded numeric representation. The positive class is recoded from original label 2; splitting and train-only preprocessing are study transformations. Attribution to the source does not imply its endorsement of these analyses.

\paragraph{ACSIncome.}
The source is the U.S. Census Bureau's ACS 2018 1-Year Public Use Microdata Sample, California person records, \texttt{csv\_pca.zip}, retrieved on 3 November 2026 \citep{uscensus2018pums}. ACSIncome is derived with folktables 0.0.12 following \citet{ding2021retiring}; the filter and weighting policy are stated in Section~\ref{sec:study}. The archive is obtained from the documented file route over HTTPS, not through the Census API. The checked Census guidance provides public-use and citation statements; no separate named license for this archive was located on the checked official pages. The MIT license for folktables software does not license the Census records. Census is the source of the original data only; these analyses and conclusions are the authors' responsibility and imply no Census endorsement.

\paragraph{Diabetes.}
We use OpenML data ID 37, version 1 \citep{openml37}. Its metadata attribute the source to the National Institute of Diabetes and Digestive and Kidney Diseases and donor Vincent Sigillito, and label the license field \emph{Public}; separate upstream license wording was not verified. This historical sample concerns women of Pima Indian heritage aged at least 21 near Phoenix. The target coding and declared zero-to-missing transformations are retained in the acquisition/loader documentation. This benchmark is not a clinical-validation cohort.

\paragraph{Source and display provenance.}
\texttt{analysis/data/manifest.json} gives retrieval URLs, dates, byte counts and SHA-256 hashes; \texttt{analysis/notes/data-attribution.md} and its source snapshots preserve the bounded source-use statements and reading provenance. Raw archives are retrieved upstream rather than redistributed. The initial pilot plot and timing table in Appendix~\ref{app:pilot} are generated together by \texttt{figures/generate\_pilot\_accounting.py}; their manifests map every panel and table entry to the unchanged pilot records. They can be regenerated from the delivery root with \texttt{python3 figures/generate\_pilot\_accounting.py --preview}. All main and background figure/table assets regenerate with \texttt{bash figures/reproduce\_main.sh} from the delivery root; per-artifact generator/input mappings are in their manifests and Appendix~\ref{app:numerical}. Main numerical inputs are the complete \texttt{results/} panel and aggregates; final figure/table commands and per-cell provenance are documented under \texttt{figures/} and \texttt{tables/}. Build the report with \texttt{bash scripts/build\_paper.sh}; this requires TeX Live with pdfLaTeX, BibTeX and latexmk, plus Poppler and ripgrep. No source in this delivery was obtained by seeking the deliberately withheld prior study; consulted materials and reading extent are recorded in \texttt{paper/source-log.md}.

\paragraph{Research-process disclosure.}
Language-model agents assisted research design, implementation, interpretation, manuscript drafting and review within separate research roles. They are part of the research process, not an estimator, reference oracle, data source or simulated participant in the empirical study. Every empirical number comes from executed numerical code and retained outputs; source/configuration fingerprints, analytic tests, saved-vector audits and independent revision-specific review constrain claims. No claim of human-subject review or clinical validation follows from this workflow.
